% Source PDF: source/普通高考/2010/2010广东文.pdf, page 1, question 11.
% Original flowchart is vector/text artwork; preserved PNG remains under img/.
% Nodes: 开始; 输入 x_1,...,x_4; s_1=0,i=1; i≤4?; s_1=s_1+x_i;
% s=(1/i)s_1; i=i+1; 输出s; 结束.
% Directed edges: 开始→输入→初始化→判定; 判定(是)→累加→平均→i更新→回主干入口→判定;
% 判定(否)→输出s→结束. The update return and output branch are independent.

\begin{tikzpicture}[
  >=Stealth,
  x=1cm,
  y=1cm,
  line width=0.45pt,
  line cap=round,
  line join=round,
  font=\normalsize,
  flowtext/.style={anchor=center, align=center,
    execute at begin node={\setlength{\parindent}{0pt}}},
  flowbox/.style={flowtext, draw, minimum height=0.68cm,
    inner xsep=4pt, inner ysep=2pt},
  terminal/.style={flowbox, rounded corners=2pt, minimum width=1.45cm},
  process/.style={flowbox, minimum width=3.25cm},
  decision/.style={flowbox, diamond, aspect=2.8, inner sep=0pt,
    minimum width=4.65cm, minimum height=1.08cm},
  io/.style={flowbox, trapezium, trapezium left angle=72,
    trapezium right angle=108, minimum height=0.78cm},
]
  \node[terminal] (start) at (0,11.60) {开始};
  \node[io, minimum width=5.20cm] (input) at (0,10.15)
    {输入 $x_1,\ldots,x_4$};
  \node[flowbox, minimum width=4.20cm] (init) at (0,8.60)
    {$s_1=0,i=1$};
  \node[decision] (test) at (0,4.55) {$i\leq 4$};
  \node[process, minimum width=4.35cm] (sum) at (5.00,4.55)
    {$s_1=s_1+x_i$};
  \node[process, minimum width=3.50cm, minimum height=1.00cm]
    (avg) at (5.00,6.25) {$s=\dfrac{1}{i}s_1$};
  \node[process, minimum width=3.20cm] (inc) at (5.00,7.95)
    {$i=i+1$};
  \node[io, minimum width=2.80cm] (output) at (0,2.65) {输出 $s$};
  \node[terminal] (finish) at (0,1.05) {结束};

  \coordinate (loopjoin) at (0,7.95);

  \draw[->] (start.south) -- (input.north);
  \draw[->] (input.south) -- (init.north);
  \draw[->] (init.south) -- (test.north);

  % 是 branch: accumulate, average, increment, then return to the main connector.
  \draw[->] (test.east) --
    node[pos=0, above left=3pt, inner sep=0pt] {是} (sum.west);
  \draw[->] (sum.north) -- (avg.south);
  \draw[->] (avg.north) -- (inc.south);
  \draw[->] (inc.west) -- (loopjoin);

  % 否 branch independently outputs s and terminates.
  \draw[->] (test.south) -- node[right=2pt, inner sep=0pt] {否} (output.north);
  \draw[->] (output.south) -- (finish.north);

  \path[use as bounding box] (-2.75,0.50) rectangle (6.05,12.00);
\end{tikzpicture}
